%======================================================================
\section{Discussion}
\label{sec:discussion}

\paragraph{Screening and estimation are different tasks.}
\enlargethispage{\baselineskip}In the same cases, the first-order estimate was the best screen we tested and a worse predictor of effect size than zero change. A screen needs only the cell types with large effects to have the largest derivatives. An estimate of the effect needs the response to stay nearly linear all the way from the intact to the silenced network, and fully silencing a cell type with a large effect often takes it well beyond that range. At half strength the screen's lead over anatomy grew while the effect-size error only matched that of zero change, and in the development models the estimate beat zero change only up to strength \DevBetterMaxStrength{}.

\paragraph{The nonlinearity is in the readout.}
Study 1 found the errors concentrated where silencing changed which directions evoke a response, and where the silenced type connects directly to the readout (\cref{tab:diagnostics}). The index is built from rectified peak responses: a direction that evokes no response contributes zero derivative, so the estimate misses responses that silencing creates and extrapolates those it removes linearly past zero. Study 2 located the cause. Keeping the network's response linear costs little: the exact index applied to the linearised trajectories beats predicting zero change, with an error ratio below 1 in \BelowOneCfRead{} of \NConfModels{} models. Linearising the index is what fails, and within the index the ratio of a vector length to a sum of responses matters most. The three single changes behave accordingly. Building the index from mean responses leaves the ratio in place and changes nothing detectable. Mean replacement shrinks the perturbation and with it the change of the ratio's denominator, which lowers the error without removing it. A second-order correction adds the ratio's curvature exactly where the ratio's Taylor series diverges, and makes the estimates far worse. \Citet{zhang2026attribution} reached a related conclusion for attribution patching in language models: most of the error comes from nonlinearities between the intervention and the output.

\paragraph{Using first-order estimates.}
For choosing which cell types to silence in a model like FlyVis, the first-order ranking is a better guide than connectivity. The sizes and signs of the chosen types' effects should then be simulated, as in the screen-then-recompute workflow of \citet{zhang2026attribution}. This advice is for silencing, which drives a type's output towards zero; for mean replacement, \citet{dhawan2026flyhack} found that estimates built for that intervention rank the sizes of its effects well, and in Study 2 the matched mean-replacement estimate beat predicting zero change at half strength. Where an estimate has to serve, clipping it to the range of the index is a minimal correction: it brought the error below that of zero change in \ClipBetterFull{} of \NConfModels{} models and left the wrong signs in place. A better correction follows from Study 2. The first-order estimate of a ratio is accurate only while the perturbation changes the denominator by a small fraction of itself, and silencing a cell type can change a motion detector's summed response by more than its own size (on the linearised trajectory, in \PostLargeDenFrac\% of cases here). For such outputs the ratio should be applied exactly to linearised inputs, or the numerator and denominator reported separately. This needs the derivatives of the \NTrajectories{} responses of each readout rather than of the index alone, and so more backward passes. Second-order corrections should be avoided for ratio outputs under large perturbations, because they inherit the divergence of the ratio's series.

The cost argument for the estimate is weak at this scale. It costs one backward pass per output, against one simulation per intervention for exact silencing, so it pays off when the interventions far outnumber the outputs, as for the many components and single output of attribution patching in language models. With \SimulatedTypes{} silencing targets and eight outputs, FlyVis is close to the break-even point, exact silencing gives the right answer, and screening a few cell types before simulating them never paid (\cref{sec:cost}).

\paragraph{The same structure in language models.}
Many outputs of language models used in patching experiments are ratios: probabilities are normalised by a softmax, and attention weights by another. The repair of attribution patching for saturated attention in \citet{kramar2024atp}, which recomputes the softmax exactly, follows the same principle as the estimator that works best here. Whether the error of attribution patching in language models comes mostly from such normalisations, or from the network between them, can be tested with the same two estimators used here, which need only one exact and one linearised forward pass per component.

\needspace{5\baselineskip}
\paragraph{Limitations.}
\begin{itemize}
  \item All \NModels{} models share one connectome, one training task and one model structure, so the intervals describe variation within this released ensemble rather than across architectures or animals.
  \item We tested moving edges at two speeds and the central T4 and T5 cells, with one family of direction-selectivity indices; other stimuli, such as natural scenes, and other outputs could give the estimate a different accuracy.
  \item Study 1's main screen is scored on cases selected after silencing, because abolished and amplified responses are excluded; the inclusive screen, which keeps every responsive readout condition and cell type, gives lower precisions with the same order of methods (\cref{sec:screening}).
  \item The index depends on peaks of rectified responses; at tied peaks the derivative depends on a tie convention, and near ties occurred in \NearTieReadouts{} of \ReadoutConditions{} readout conditions.
  \item Study 2 was designed after Study 1's results were known. Its arms~(b) and~(c), the analysis of where the nonlinearity lies and the flag rule were computed from Study 1's saved simulations of the confirmation models, which Study 1 had already analysed; they were prespecified but are not blind confirmations. The split of the readout into its two stages and a second set of predictions were added after Study 2's development results were read, before its freeze. The explanation of the second-order failure in \cref{sec:second} is post hoc.
  \item The second-order term is the one automatic differentiation gives, which treats every rectifier as linear; a finite-difference term would include threshold crossings at a chosen step size. Integrated gradients along each cell type's own silencing path~\citep{sundararajan2017axiomatic} converge to the exact change, but each cell type then needs several simulations and backward passes, more than the single simulation they would replace. Mean replacement used each cell's average over the stimulus set; other averages are possible.
  \item Silencing in the model scales a type's synaptic output for the duration of one stimulus (at most \WindowSlow{}~s); genetic silencing in flies acts for longer and through other mechanisms, so the biological comparisons in \cref{app:biology} are qualitative.
\end{itemize}

%======================================================================
\section{Conclusion}
\label{sec:conclusion}

In the \NModels{} FlyVis models, first-order estimates of cell-type silencing are a good screen: they pick out the cell types whose silencing changes the direction selectivity of T4 and T5 more reliably than any ranking from connectivity, trained weights or input currents that we tested. They are a poor estimate of the effect itself. At full silencing their errors exceed those of predicting zero change, a fifth of large effects get the wrong sign, and among large effects their ranking is no better than that of anatomy. The cause is the readout, not the network: applying the index exactly to the linearised trajectories beats predicting zero change, and within the index the ratio matters most. A second-order correction makes the estimates worse, and mean replacement, a smaller perturbation, lowers the error without removing it. At this scale exact silencing costs about as much as the estimate, and when the size or sign of the effect of driving a cell type's output towards zero matters, it should be simulated.

\paragraph{Code and data availability.}
\enlargethispage{\baselineskip}The code, frozen protocols and records of both studies, the saved outputs from which every number, table and figure in this paper is computed, and the paper's source are available at \url{https://github.com/Jeevash23/interp-flyvis-first-order}. The simulated voltage traces are not included because of their size; the released code computes them again from the public FlyVis models. All of Study 1's saved responses were rechecked from their traces, one silencing case per model was replayed with the original FlyVis implementation, and Study 2 checked every confirmation model against Study 1's saved records before using its outcomes (\cref{app:verification}).
